\documentclass[10pt, aspectratio=169]{beamer}
% Preámbulo modular para presentaciones Beamer (Modelación Estadística)
% Diseño y paleta institucional del Tecnológico de Monterrey

\documentclass[aspectratio=169,xcolor={dvipsnames,table}]{beamer}

\usepackage[utf8]{inputenc}
\usepackage[spanish,mexico]{babel}
\usepackage{amsmath,amssymb,amsfonts,mathtools}
\usepackage{graphicx}
\usepackage{listings}
\usepackage{booktabs}
\usepackage{tikz}
\usetikzlibrary{matrix,arrows,decorations.pathmorphing,shapes.geometric,positioning}

% Paleta Institucional Tecnológico de Monterrey
\definecolor{TecAzul}{HTML}{0039A6}       % Azul TEC: color primario institucional
\definecolor{TecAzulOscuro}{HTML}{1A2E51} % Azul marino oscuro
\definecolor{TecRojo}{HTML}{EC2661}       % Rojo de acento (paleta miTec)
\definecolor{TecGris}{HTML}{646464}       % Gris neutro para comentarios y subtítulos
\definecolor{TecFondoBloque}{HTML}{F4F6F9}% Fondo suave para bloques

% Configuración de Tema Beamer
\usetheme{Boadilla}
\usecolortheme[named=TecAzulOscuro]{structure}
\setbeamercolor{alerted text}{fg=TecRojo}
\setbeamercolor{palette primary}{bg=TecAzulOscuro,fg=white}
\setbeamercolor{palette secondary}{bg=TecAzul,fg=white}
\setbeamercolor{palette tertiary}{bg=TecAzulOscuro!80!black,fg=white}
\setbeamercolor{title}{fg=TecAzulOscuro,bg=TecFondoBloque}
\setbeamercolor{frametitle}{fg=TecAzulOscuro,bg=TecFondoBloque}

% Bloques personalizados elegantes
\setbeamercolor{block title}{bg=TecAzulOscuro,fg=white}
\setbeamercolor{block body}{bg=TecFondoBloque,fg=black}
\setbeamercolor{block title alerted}{bg=TecRojo,fg=white}
\setbeamercolor{block body alerted}{bg=TecRojo!8,fg=black}
\setbeamercolor{block title example}{bg=TecAzul,fg=white}
\setbeamercolor{block body example}{bg=TecAzul!8,fg=black}

% Redondear bordes y añadir sombra sutil
\setbeamertemplate{blocks}[rounded][shadow=true]
\setbeamertemplate{navigation symbols}{}

% Pie de página limpio con número de diapositiva
\setbeamertemplate{footline}{
  \leavevmode%
  \hbox{%
  \begin{beamercolorbox}[wd=.7\paperwidth,ht=2.25ex,dp=1ex,left]{author in head/foot}%
    \hspace*{2ex}\insertshortauthor~~(\insertshortinstitute) \hfill \insertshorttitle
  \end{beamercolorbox}%
  \begin{beamercolorbox}[wd=.3\paperwidth,ht=2.25ex,dp=1ex,right]{date in head/foot}%
    \insertshortdate{}\hspace*{2em}
    \insertframenumber{} / \inserttotalframenumber\hspace*{2ex}
  \end{beamercolorbox}}%
  \vskip0pt%
}

% Configuración de Listings de Python para visualización pedagógica de código
\lstset{
    language=Python,
    basicstyle=\ttfamily\footnotesize,
    keywordstyle=\color{TecAzulOscuro}\bfseries,
    stringstyle=\color{TecRojo},
    commentstyle=\color{TecGris}\itshape,
    showstringspaces=false,
    breaklines=true,
    frame=lines,
    rulecolor=\color{TecAzulOscuro!30},
    backgroundcolor=\color{TecFondoBloque},
    aboveskip=0.5em,
    belowskip=0.5em,
    literate={á}{{\'a}}1 {é}{{\'e}}1 {í}{{\'i}}1 {ó}{{\'o}}1 {ú}{{\'u}}1
             {Á}{{\'A}}1 {É}{{\'E}}1 {Í}{{\'I}}1 {Ó}{{\'O}}1 {Ú}{{\'U}}1
             {ñ}{{\~n}}1 {Ñ}{{\~N}}1 {¿}{{?`}}1 {¡}{{!`}}1
}

% Metadatos institucionales por defecto
\title[Modelación Estadística]{Modelación Estadística para Ciencia de Datos e Ingeniería}
\author[J. Castillo Colmenares]{Juliho David Castillo Colmenares}
\institute[Tec de Monterrey]{Tecnológico de Monterrey \\ Escuela de Ingeniería y Ciencias}
\date{\today}

% Comandos y macros estadísticas compartidas para las presentaciones Beamer (Metropolis)

% Conjuntos numéricos básicos
\newcommand{\R}{\mathbb{R}}
\newcommand{\N}{\mathbb{N}}
\newcommand{\Q}{\mathbb{Q}}
\newcommand{\Z}{\mathbb{Z}}
\newcommand{\set}[1]{\left\{ #1 \right\}}
\newcommand{\abs}[1]{\left|#1\right|}
\newcommand{\norm}[1]{\left\| #1 \right\|}

% Comandos estadísticos y probabilísticos del curso
\newcommand{\Var}{\operatorname{Var}}
\newcommand{\cov}{\operatorname{Cov}}
\newcommand{\corr}{\rho}
\newcommand{\comb}[2]{\begin{pmatrix} #1 \\ #2 \end{pmatrix}}
\newcommand{\s}{\sigma}
\newcommand{\particion}{\mathfrak{P}}
\newcommand{\card}[1]{\#\left( #1 \right)}
\newcommand{\seq}[3]{\set{#1}_{#2}^{#3}}
\newcommand{\E}{\mathbb{E}}
\newcommand{\Prob}{\mathbb{P}}

% Entornos pedagógicos y cajas en español académico natural
\newenvironment{discusion}{%
  \begin{alertblock}{Pregunta de Discusión}%
}{%
  \end{alertblock}%
}

\newenvironment{reflexion}{%
  \begin{alertblock}{Reflexión para el Aula}%
}{%
  \end{alertblock}%
}

\newenvironment{paradoja}{%
  \begin{alertblock}{Paradoja de Exploración}%
}{%
  \end{alertblock}%
}

\newenvironment{motivacion}{%
  \begin{exampleblock}{Motivación: Ciencia de Datos en la Práctica}%
}{%
  \end{exampleblock}%
}

\newenvironment{retoclase}{%
  \begin{block}{Reto Rápido en Equipo}%
}{%
  \end{block}%
}


\title[Calidad de Estimadores]{Sección 06.01: Estimación Puntual, Insesgadez, Eficiencia y Consistencia}
\subtitle{Sesgo, Error Cuadrático Medio, y la Cota de Cramér-Rao}
\author[J. Castillo Colmenares]{Juliho Castillo Colmenares}
\institute{Tecnológico de Monterrey}
\date{\vspace{-1.2cm}}

\begin{document}

% Slide 1: Portada
\begin{frame}[plain]
  \titlepage
\end{frame}

% Slide 2: Hoja de Ruta
\begin{frame}{Hoja de Ruta --- Capítulo 06: Estimación y su Relación con Ciencia de Datos}
  \footnotesize
  ¿Dónde nos encontramos en el desarrollo modular del capítulo? \pause
  \begin{itemize}\itemsep=0.08em
    \item \textbf{06.01} Estimación Puntual, Insesgadez, Eficiencia y Consistencia \emph{(Hoy)} \pause
    \item \textbf{06.02} Método de Momentos (MoM) \pause
    \item \textbf{06.03} Estimación por Máxima Verosimilitud (MLE) y Score \pause
    \item \textbf{06.04} Intervalos de Confianza para Medias Poblacionales ($Z$ y $t$) \pause
    \item \textbf{06.05} Intervalos de Confianza para Varianzas ($\chi^2$) y Proporciones
  \end{itemize}
  \vspace{-0.05cm}
  \begin{block}{Objetivo de la Sesión}
    Formalizar los criterios que definen a un ``buen'' estimador --- sesgo, error cuadrático medio, eficiencia relativa y consistencia --- y establecer la Cota de Cramér-Rao como el límite teórico de precisión que ningún estimador insesgado puede superar.
  \end{block}
\end{frame}

% Slide 3: Motivación
\begin{frame}{De las Distribuciones de Muestreo a la Teoría de Estimación}
  \footnotesize
  \begin{motivacion}
    En el Capítulo 05 estudiamos la distribución muestral de $\bar X$ y $S^2$. Ahora damos un paso atrás: ¿qué hace que un estimador sea ``bueno''? Antes de construir estimadores (Secciones 06.02-06.03), necesitamos un marco formal para \emph{evaluarlos}.
  \end{motivacion}

  \vspace{0.15cm}
  \pause
  \begin{itemize}\itemsep=0.1cm
    \item \textbf{Insesgadez:} ¿el estimador acierta ``en promedio''? \pause
    \item \textbf{Eficiencia:} entre varios estimadores insesgados, ¿cuál tiene menor varianza? \pause
    \item \textbf{Consistencia:} ¿mejora el estimador conforme crece la muestra?
  \end{itemize}
\end{frame}

% Slide 4: Sesgo y ECM
\begin{frame}{Sesgo y Error Cuadrático Medio}
  \footnotesize
  Sea $\hat\theta$ un estimador puntual de $\theta$. \pause

  \vspace{0.1cm}
  \begin{block}{Definiciones}
    $$\text{Sesgo}(\hat\theta) = E(\hat\theta)-\theta, \qquad \text{ECM}(\hat\theta) = E\left[(\hat\theta-\theta)^2\right]$$
  \end{block}

  \vspace{0.1cm}
  \pause
  \begin{alertblock}{Descomposición Sesgo-Varianza}
    $$\text{ECM}(\hat\theta) = \Var(\hat\theta) + \left[\text{Sesgo}(\hat\theta)\right]^2$$
    Un estimador ligeramente sesgado puede tener menor ECM que uno insesgado si reduce suficientemente la varianza --- la base del \emph{trade-off sesgo-varianza} en Ridge/Lasso.
  \end{alertblock}
\end{frame}

% Slide 5: Eficiencia y Cramér-Rao
\begin{frame}{Eficiencia Relativa y la Cota de Cramér-Rao}
  \footnotesize
  Para comparar estimadores insesgados: \pause

  \vspace{0.1cm}
  \begin{block}{Eficiencia Relativa}
    $$\text{Ef}(\hat\theta_1,\hat\theta_2) = \dfrac{\Var(\hat\theta_2)}{\Var(\hat\theta_1)}$$
  \end{block}

  \vspace{0.1cm}
  \pause
  \begin{alertblock}{Cota Inferior de Cramér-Rao (CRLB)}
    $$\Var(\hat\theta) \geq \dfrac{1}{n\,I(\theta)}, \qquad I(\theta) = -E\left[\dfrac{\partial^2}{\partial\theta^2}\ln f(X\mid\theta)\right]$$
    Un estimador que alcanza esta cota es \textbf{UMVUE} (Estimador Insesgado de Varianza Mínima Uniformemente): el mejor posible.
  \end{alertblock}
\end{frame}

% Slide 6: Consistencia
\begin{frame}{Consistencia}
  \footnotesize
  Una propiedad asintótica: ¿mejora el estimador con más datos? \pause

  \vspace{0.1cm}
  \begin{block}{Definición}
    $$\hat\theta_n \xrightarrow{P} \theta \iff \forall\varepsilon>0,\ \lim_{n\to\infty} P(|\hat\theta_n-\theta|\ge\varepsilon) = 0$$
  \end{block}

  \vspace{0.1cm}
  \pause
  \begin{alertblock}{Condición Suficiente Práctica}
    Por Chebyshev, $P(|\hat\theta_n-\theta|\ge\varepsilon) \le \text{ECM}(\hat\theta_n)/\varepsilon^2$. Si $\text{ECM}(\hat\theta_n)\to 0$ (sesgo y varianza $\to 0$), entonces $\hat\theta_n$ es consistente. Esta es la vía más común para demostrar consistencia.
  \end{alertblock}
\end{frame}

% Slide 7: Aplicaciones
\begin{frame}{Aplicaciones en Ciencia de Datos}
  \footnotesize
  Estos criterios son el lenguaje común de toda la estadística inferencial y el aprendizaje automático: \pause

  \vspace{0.15cm}
  \begin{itemize}\itemsep=0.12cm
    \item \textbf{Regularización (Ridge/Lasso):} acepta sesgo controlado a cambio de menor varianza y mejor error de generalización. \pause
    \item \textbf{Validación cruzada:} estima el ECM de un modelo fuera de muestra. \pause
    \item \textbf{Selección de estimadores:} la eficiencia relativa guía la elección entre métodos (MoM vs. MLE, Secciones 06.02-06.03). \pause
    \item \textbf{Big data:} la consistencia garantiza que más datos siempre ayudan, incluso con modelos aproximados.
  \end{itemize}
\end{frame}

% Slide 8: Lab Python (Bloque 1)
\begin{frame}[fragile]{Laboratorio en Python: Sesgo, ECM y Eficiencia}
  \scriptsize
  Descomposición sesgo-varianza y eficiencia relativa comparando $\bar X$ contra un estimador ingenuo (\texttt{06.01\_point\_estimation\_quality.py}):
  {\lstset{basicstyle=\fontsize{4pt}{4.8pt}\selectfont\ttfamily,aboveskip=0.1em,belowskip=0.1em}
  \lstinputlisting[firstline=17, lastline=36]{../../code/06_estimacion_estadistica/06.01_point_estimation_quality.py}}
\end{frame}

% Slide 9: Lab Python (Bloque 2)
\begin{frame}[fragile]{Laboratorio en Python: Estimador de Encogimiento}
  \scriptsize
  Cálculo del factor de encogimiento óptimo $c^*$ que minimiza el ECM, verificado por búsqueda en malla:
  {\lstset{basicstyle=\fontsize{4pt}{4.8pt}\selectfont\ttfamily,aboveskip=0.1em,belowskip=0.1em}
  \lstinputlisting[firstline=39, lastline=59]{../../code/06_estimacion_estadistica/06.01_point_estimation_quality.py}}
\end{frame}

% Slide 10: Lab Python (Bloque 3)
\begin{frame}[fragile]{Laboratorio en Python: Consistencia vía Chebyshev}
  \scriptsize
  Cota de Chebyshev vs. probabilidad empírica Monte Carlo para la proporción muestral:
  {\lstset{basicstyle=\fontsize{4pt}{4.8pt}\selectfont\ttfamily,aboveskip=0.1em,belowskip=0.1em}
  \lstinputlisting[firstline=62, lastline=82]{../../code/06_estimacion_estadistica/06.01_point_estimation_quality.py}}
\end{frame}

% Slide 11: Salida Terminal
\begin{frame}[fragile]{Laboratorio en Python: Salida en Terminal}
  \scriptsize
  Salida estándar generada al ejecutar el script del laboratorio en Python:
  \vspace{0.02cm}
  \begin{block}{Salida Estándar en Consola}
    \fontsize{4pt}{5pt}\selectfont
    \begin{verbatim}
=== Block 1: Bias-Variance-MSE Decomposition and Relative Efficiency ===
Xbar: bias=0.009, Var=6.24, MSE=6.24 | T1: bias=0.023, Var=100.22, MSE=100.22
Relative efficiency Ef(Xbar, T1) = 16.07 (theoretical: n=16)

=== Block 2: Shrinkage Estimator ===
c*=0.6429 | MSE(c=1)=5.00 | MSE(c*)=3.21 | reduction: 35.7%

=== Block 3: Consistency via Chebyshev ===
n=100: bound=1.00, empirical=0.319 | n=1000: bound=0.10, empirical=0.001
Minimum n for bound<=0.10: n>=1000
    \end{verbatim}
  \end{block}
\end{frame}

% Slide 12A: Ejercicio Nivel 1 (Enunciado)
\begin{frame}{Ejercicio en Clase --- Nivel 1: Fundamental (1/2)}
  \footnotesize
  \begin{block}{Problema 6.1.1 --- Insesgadez de Combinaciones Lineales (Enunciado)}
    Sean $X_1, X_2$ iid con $E(X_i)=\mu$. Considere $\hat\theta_w = 0.3X_1+0.7X_2$. Demuestre que $\hat\theta_w$ es insesgado para $\mu$, y generalice la condición sobre los pesos.
  \end{block}

  \vspace{0.15cm}
  \pause
  \begin{alertblock}{Planteamiento}
    \begin{itemize}\itemsep=0.04cm
      \item Use linealidad de la esperanza: $E(\hat\theta_w)=0.3E(X_1)+0.7E(X_2)$.
    \end{itemize}
  \end{alertblock}
\end{frame}

% Slide 12B: Ejercicio Nivel 1 (Resolución)
\begin{frame}{Ejercicio en Clase --- Nivel 1: Fundamental (2/2)}
  \scriptsize
  Aplicamos linealidad: \pause

  \vspace{0.05cm}
  \begin{block}{Cálculo}
    $$E(\hat\theta_w) = 0.3\mu + 0.7\mu = (0.3+0.7)\mu = \mu$$
  \end{block}

  \vspace{0.05cm}
  \pause
  \begin{alertblock}{Generalización}
    \scriptsize
    Para $\hat\theta=\sum_i w_i X_i$ con $E(X_i)=\mu$ para todo $i$: $E(\hat\theta)=\left(\sum_i w_i\right)\mu$. El estimador es insesgado si y solo si $\sum_i w_i = 1$ --- la media muestral ($w_i=1/n$) es solo un caso particular de esta familia.
  \end{alertblock}
\end{frame}

% Slide 13A: Ejercicio Nivel 2 (Enunciado)
\begin{frame}{Ejercicio en Clase --- Nivel 2: Operativo (1/2)}
  \footnotesize
  \begin{block}{Problema 6.1.6 --- Consistencia vía Chebyshev (Enunciado)}
    Un estimador insesgado $\hat\theta_n$ tiene $\Var(\hat\theta_n)=16/n$. Use Chebyshev para acotar $P(|\hat\theta_n-\theta|\ge 1)$ cuando $n=64$, y explique por qué esto confirma consistencia.
  \end{block}

  \vspace{0.15cm}
  \pause
  \begin{alertblock}{Estrategia}
    \scriptsize
    $P(|\hat\theta_n-\theta|\ge\varepsilon) \le \Var(\hat\theta_n)/\varepsilon^2$.
  \end{alertblock}
\end{frame}

% Slide 13B: Ejercicio Nivel 2 (Resolución)
\begin{frame}{Ejercicio en Clase --- Nivel 2: Operativo (2/2)}
  \scriptsize
  Sustituimos en la cota de Chebyshev: \pause

  \vspace{0.05cm}
  \begin{block}{Cálculo}
    $$P(|\hat\theta_{64}-\theta|\ge 1) \le \frac{16/64}{1^2} = 0.25$$
  \end{block}

  \vspace{0.05cm}
  \pause
  \begin{alertblock}{Interpretación}
    \scriptsize
    Como $16/n \to 0$ cuando $n\to\infty$ para cualquier $\varepsilon$ fijo, la cota de Chebyshev tiende a $0$, confirmando que $\hat\theta_n \xrightarrow{P} \theta$: el estimador es consistente.
  \end{alertblock}
\end{frame}

% Slide 14A: Ejercicio Nivel 3 (Enunciado)
\begin{frame}{Ejercicio en Clase --- Nivel 3: Analítico (1/2)}
  \footnotesize
  \begin{block}{Problema 6.1.7 --- CRLB para la Media Normal (Enunciado)}
    Sea $X_1,\ldots,X_n \sim N(\mu,\sigma^2)$ con $\sigma^2$ conocida. Calcule $I(\mu)$, derive la CRLB para $\Var(\hat\mu)$, y demuestre que $\bar X$ la alcanza exactamente.
  \end{block}

  \vspace{0.1cm}
  \pause
  \begin{alertblock}{Estrategia}
    \scriptsize
    Derive $\ln f(x\mid\mu)$ dos veces respecto a $\mu$ para obtener $I(\mu)$.
  \end{alertblock}
\end{frame}

% Slide 14B: Ejercicio Nivel 3 (Resolución)
\begin{frame}{Ejercicio en Clase --- Nivel 3: Analítico (2/2)}
  \scriptsize
  Calculamos la Información de Fisher: \pause

  \vspace{0.05cm}
  \begin{block}{Derivación}
    \scriptsize
    \begin{align*}
      \frac{\partial}{\partial\mu}\ln f = \frac{x-\mu}{\sigma^2}, \qquad \frac{\partial^2}{\partial\mu^2}\ln f = -\frac{1}{\sigma^2} \implies I(\mu) = \frac{1}{\sigma^2}.
    \end{align*}
    \begin{align*}
      \text{CRLB} = \frac{1}{nI(\mu)} = \frac{\sigma^2}{n} = \Var(\bar X). \quad \qedhere
    \end{align*}
  \end{block}

  \vspace{0.05cm}
  \pause
  \begin{alertblock}{Interpretación}
    \scriptsize
    Como $\bar X$ alcanza la CRLB exactamente, es el UMVUE para $\mu$: ningún estimador insesgado puede tener menor varianza.
  \end{alertblock}
\end{frame}

% Slide 15A: Ejercicio Nivel 4 (Enunciado)
\begin{frame}{Ejercicio en Clase --- Nivel 4: Desafiante (1/2)}
  \footnotesize
  \begin{block}{Problema 6.1.9 --- Estimador de Encogimiento Óptimo (Enunciado)}
    Con $\Var(\hat\theta)=v=5$ y $\theta=3$, evalúe numéricamente el factor óptimo $c^*=\theta^2/(v+\theta^2)$ y compare $\text{ECM}(\hat\theta_{c^*})$ contra $\text{ECM}(\hat\theta)=v$ (caso insesgado).
  \end{block}

  \vspace{0.15cm}
  \pause
  \begin{alertblock}{Estrategia}
    \scriptsize
    Sustituya $v=5$, $\theta=3$ en $c^*$ y luego en $\text{ECM}(c)=c^2v+(c-1)^2\theta^2$.
  \end{alertblock}
\end{frame}

% Slide 15B: Ejercicio Nivel 4 (Resolución)
\begin{frame}{Ejercicio en Clase --- Nivel 4: Desafiante (2/2)}
  \scriptsize
  Sustituimos numéricamente: \pause

  \vspace{0.05cm}
  \begin{block}{Cálculo}
    \scriptsize
    \begin{align*}
      c^* = \frac{9}{5+9} = \frac{9}{14} \approx 0.643, \qquad \text{ECM}(c^*) \approx 2.066+1.148 = 3.214.
    \end{align*}
  \end{block}

  \vspace{0.05cm}
  \pause
  \begin{alertblock}{Interpretación}
    \scriptsize
    Comparado con $\text{ECM}(c=1)=5$, el ECM se reduce aproximadamente un $35.7\%$ al aceptar un sesgo controlado: un ejemplo concreto de por qué la regularización (Ridge/Lasso) sacrifica insesgadez por menor error total.
  \end{alertblock}
\end{frame}

% Slide 16: Cierre
\begin{frame}{Cierre de Sección: Criterios de Calidad de Estimadores}
  \begin{reflexion}
    Antes de aprender a construir estimadores (MoM, MLE), necesitábamos el lenguaje para evaluarlos. Sesgo, ECM, eficiencia y consistencia son los cuatro pilares que usaremos --- explícita o implícitamente --- en cada método del resto del capítulo.
  \end{reflexion}

  \vspace{0.2cm}
  \pause
  \begin{block}{Lecciones Clave}
    \begin{itemize}\itemsep=0.08cm
      \item \textbf{ECM = Varianza + Sesgo$^2$:} un pequeño sesgo puede reducir el error total.
      \item \textbf{Cramér-Rao:} $\Var(\hat\theta) \ge 1/(nI(\theta))$; alcanzarla define al UMVUE.
      \item \textbf{Consistencia:} si $\text{ECM}(\hat\theta_n)\to 0$, entonces $\hat\theta_n \xrightarrow{P} \theta$.
    \end{itemize}
  \end{block}
\end{frame}

% Slide 17: Perspectiva Modular
\begin{frame}{Perspectiva Modular --- El Próximo Reto}
  \scriptsize
  \begin{retoclase}
    Con los criterios de calidad establecidos, la Sección 06.02 introducirá el Método de Momentos: un procedimiento sistemático para \emph{construir} estimadores igualando momentos muestrales y poblacionales. La Sección 06.03 hará lo propio con el Método de Máxima Verosimilitud, el método más poderoso y ampliamente usado en ciencia de datos moderna.
  \end{retoclase}

  \vspace{0.08cm}
  \pause
  \begin{alertblock}{Preparación Recomendada}
    \begin{itemize}\itemsep=0.06cm
      \item Repasa la definición de momentos poblacionales y muestrales (Capítulo 03/04).
      \item Reflexiona sobre por qué distintos métodos de construcción pueden producir el mismo estimador (como ocurre con la Exponencial) o estimadores distintos (como con la Uniforme).
    \end{itemize}
  \end{alertblock}
\end{frame}

\end{document}
